\scanpage{161}
Supposons que pour tout $X\in\mathcal V$, \inserted{on ait un}
$(\uline h(X)=(\uline h^i(X)))$
\deleted{\unsure{soit muni d'un}}
un ensemble $\Phi_X$ de formes fondamentales
$\{\varphi\}$, chaque $\varphi=(\varphi_i)$,
$\varphi_i\in\operatorname{Hom}(\uline h^i(X)\otimes\uline h^i(X),
\Lambda(-i))$, de façon que les conditions suivantes
soient vérifiées :

(i) Si $\varphi\in\Phi_X$, $\psi\in\Phi_Y$, alors
$\varphi\oplus\psi\in\Phi_{X\amalg Y}$.

(ii) Si $\varphi\in\Phi_X$, $\psi\in\Phi_Y$,
\deleted{et $u:X\longrightarrow Y$, d'où
$h(u):\uline h^i(Y)\longrightarrow\uline h^i(X)$, \unsure{pourrions}}
alors $\varphi\otimes\psi\in\Phi_{X\unsure{\times|\otimes}Y}$

$(\mathrm{iii}_0)$ Sur
\deleted{$\varphi\in\Phi_X$, alors $\varphi(0)$ sur $\unsure{X|\overline X}=e$, donc
$\uline h^0(e)\unsure{\otimes|\simeq}\uline h^0(e)\simeq Q$, \unsure{tous}}
\inserted{$h^0(X)$, \unsure{considéré de} \unsure{dim} finie sur $Q$, les $\varphi\in\Phi_X$ sont}
des formes \uline{définies positives}
\margin{[il suffit $X=e$ \deleted{\unsure{lorsque}}]}

\unsure{Sans doute on le supprimera} \unsure{par}
\deleted{(iii)}

(iii) Si $\varphi\in\Phi_X$, $\psi\in\Phi_Y$,
$u:X\mathbin{-}Y$ d'où
$u^{(i)}:h^i(\unsure{Y|X})\mathbin{-}h^i(Y)$,
\unsure{prenant} $\unsure{\delta|s}u^{(i)}:h^i(Y)\mathbin{-}h^i(X)$,
[\unsure{relativ.}: $\varphi,\psi$], on a
\[
\boxed{\operatorname{Tr}(\unsure{\delta|s}u^{(i)}u^i)\geq0,
\quad\text{égalité ssi }u^{(i)}=0}
\]
\mspage{160}
